\documentclass[10pt]{article}
\usepackage[letterpaper,margin=0.58in]{geometry}
\usepackage{booktabs}
\usepackage{tabularx}
\usepackage{xcolor}
\usepackage{microtype}
\usepackage{hyperref}
\usepackage{enumitem}
\setlist{nosep,leftmargin=*}
\setlength{\parindent}{0pt}
\setlength{\parskip}{3pt}
\newcommand{\todo}[1]{\textcolor{red}{[REPLACE: #1]}}

\title{\vspace{-1.3em}ITCS 6169/8169 Assignment 1: The CNN Challenge}
\author{\todo{Name}}
\date{}

\begin{document}
\maketitle
\vspace{-2.1em}

\textbf{Code and checkpoint.}
\href{https://github.com/REPLACE/REPLACE}{GitHub repository}\todo{replace URL}.
The repository contains
the exact configuration, split, code, environment record, and checkpoint used
for the result below.

\section*{Final result and validation strategy}
The selected model achieved \textbf{\todo{test accuracy}} on the labeled test
set. It was selected using a fixed, class-stratified 80/20 split of the 2,400
provided training images (1,920 train and 480 validation images; seed 0). The
selected checkpoint obtained \textbf{\todo{validation accuracy}} validation
accuracy at epoch \todo{best epoch}. The test set was not used for model or
hyperparameter selection.

\section*{Final recipe}
The final classifier was \todo{CNN architecture} initialized from
\todo{pretraining dataset or ``random initialization''}. Images were processed
at \todo{resolution} using \todo{augmentation summary}. I trained with
\todo{optimizer}, batch size \todo{batch size}, initial backbone/head learning
rates \todo{learning rates}, \todo{schedule}, weight decay \todo{weight decay},
and label smoothing \todo{value}. \todo{Explain which one or two choices mattered
most, using validation evidence rather than listing settings.}

\section*{Experimental journey}
\begin{center}
\small
\begin{tabularx}{\linewidth}{@{}lXrr@{}}
\toprule
Experiment & Controlled change and observation & Val. acc. & Time \\
\midrule
Starter TNet & Grayscale $64\times64$ sanity-check baseline. & \todo{--} & \todo{--} \\
Linear probe & ImageNet ResNet-18; classifier trained with frozen backbone. & \todo{--} & \todo{--} \\
Fine-tuning & Unfroze backbone after three epochs; basic augmentation. & \todo{--} & \todo{--} \\
Augmentation & Changed only the augmentation profile from basic to strong. & \todo{--} & \todo{--} \\
Final model & \todo{Final architecture/change and concise observation.} & \todo{--} & \todo{--} \\
\bottomrule
\end{tabularx}
\end{center}

\section*{Failure analysis}
I expected \todo{failed experiment} to help because \todo{hypothesis}. Instead,
validation accuracy \todo{what happened}, while the learning curves showed
\todo{diagnostic evidence}. This suggests \todo{lesson and resulting decision}.
This experiment was not selected based on its test-set behavior.

\section*{Error analysis}
The final confusion matrix showed that \todo{difficult/confused classes and
counts}. Visual inspection suggests \todo{plausible image-level reason}. A
remaining limitation is \todo{limitation}, which could be tested in future work
with \todo{controlled follow-up}.

\section*{AI + human}
\todo{One example where an AI coding assistant helped and one where my
judgment was needed.}

\section*{Reproducibility}
The final run can be reproduced with:
\begin{verbatim}
python train.py --config configs/FINAL_CONFIG.yaml
python evaluate.py --checkpoint runs/FINAL/best.pt \
  --data-root data --output-dir runs/FINAL/test_evaluation
\end{verbatim}
Software versions, seed, class order, resolved configuration, and split indices
are stored with the run artifacts.

\newpage
\section*{References}
\small
K. He et al., ``Deep Residual Learning for Image Recognition,'' CVPR, 2016.

M. Tan and Q. Le, ``EfficientNet: Rethinking Model Scaling for Convolutional
Neural Networks,'' ICML, 2019.

PyTorch and torchvision documentation, versions recorded in the repository.

\end{document}
